\section{Coupling constraints}\label{sec:constraints}

The corrected-grid bound of Proposition~\ref{prop:cellwise} follows from
rounding each coordinate independently to an endpoint of its grid interval,
and every rounded point is feasible because the domain is a product. A coupling constraint, even one as sparse as $x_1=x_2$, destroys
this: independent rounding leaves the feasible set, and the minimum over
feasible grid points need not be a lower bound. This section shows that
totally unimodular (TU) coupling of the continuous variables, with data
aligned to a uniform mesh, admits a \emph{feasible} mean-preserving rounding.
Filtering, certificates, accuracy-independent node counts and exact output
then carry over. The price is visible in two places. The rounding is
correlated, so a bound on the full continuous Hessian replaces coordinate
curvature; Example~\ref{ex:tu-fullcurv} shows that this is necessary. The
analysis uses a uniform aligned mesh, so the number of nodes per coordinate
is of order $\sqrt{n_c\kappa_c}$ instead of logarithmic in $n$;
Section~\ref{sec:tu-limits} shows that the algorithm incurs this cost, and
that two natural non-uniform alternatives, misaligned coordinate grids and a
common graded pattern, give invalid bounds with curvature-only corrections.
Section~\ref{sec:tu-complexity} states when the algorithm is polynomial and
why it is not fixed-parameter tractable.

\subsection{Model}\label{sec:tu-model}

\begin{definition}[Aligned TU model]\label{def:tu-model}
Let $n_c\ge1$ and $n_d\ge0$. The feasible set is
\[
X=\bigl\{(x,z)\in\R^{n_c}\times\Z^{n_d}:\ \ell\le x\le u,\ z\in\mathcal Z,\
Ax+Bz\le b\bigr\},\qquad \mathcal Z=Z_1\times\dots\times Z_{n_d},
\]
with the following data.
\begin{enumerate}[label=(\roman*)]
\item $A\in\{0,\pm1\}^{m\times n_c}$ is totally unimodular; $B$ and $b$ are
rational.
\item $\ell<u$ componentwise (fixed continuous coordinates have been
substituted), and every $Z_k$ is a nonempty, explicitly listed finite set of
integers.
\item A rational \emph{mesh unit} $\eta>0$ is given such that
\begin{equation}\label{eq:tu-align}
\ell/\eta\in\Z^{n_c},\qquad u/\eta\in\Z^{n_c},\qquad
(b-Bz)/\eta\in\Z^m\quad\text{for every }z\in\mathcal Z .
\end{equation}
For integral $\ell,u,B,b$ one may take $\eta=1$. Write $\Delta_\eta$ for the
denominator of $\eta$ in lowest terms.
\item $F(x,z)=\sum_{a\in\mathcal A}f_a(v_{S_a})$ with $v=(x,z)$, and a tree
decomposition with $N$ nodes and bags of size at most $p$ is given such that
every scope $S_a$ and every row support
$\{i:A_{ri}\ne0\}\cup\{k:B_{rk}\ne0\}$ is contained in some bag.
\item (Directional curvature.) A linear subspace $\mathcal W\subseteq\R^{n_c}$
and a rational $\bar L>0$ are given such that, for every $z\in\mathcal Z$,
$F(\cdot,z)$ is twice continuously differentiable on an open set containing
$[\ell,u]$ and
\begin{equation}\label{eq:tu-curv}
w^{\T}\nabla^2_{xx}F(x,z)\,w\le\bar L\norm w^2\qquad
(x\in[\ell,u],\ w\in\mathcal W).
\end{equation}
Either $\mathcal W=\R^{n_c}$, or $\mathcal W=\ker C$, where the rows of $C$
are the continuous parts $A_r$ of some \emph{equality pairs}, that is, pairs
of rows $r,r'$ with $(A_{r'},B_{r'},b_{r'})=-(A_r,B_r,b_r)$.
\end{enumerate}
\end{definition}

Equalities are thus written as two opposite inequalities; this preserves
total unimodularity. Only the continuous columns $A$ must be TU: the integer
columns $B$ are arbitrary, because the rounding below never changes $z$.
Throughout we write $n=n_c+n_d$,
\[
s=\max_{i\le n_c}(u_i-\ell_i),\quad K_Z=\max\Bigl\{1,\max_k\abs{Z_k}\Bigr\},\quad
h_j=\eta\,2^{-j},\quad E_j=\frac{n_c\bar L h_j^2}{8},
\]
$\mathcal S=\argmin_XF$, $\OPT=\min_XF$, and
$\mathcal S_i=\{v_i:v\in\mathcal S\}$ for the projection of $\mathcal S$ on
coordinate $i$ (continuous or discrete). Since the sets $Z_k$ are listed,
$K_Z\le I$. As in CT with a common mesh (Lemma~\ref{lem:commonmesh}), $h_j$
is the same for all continuous coordinates; its base is the mesh unit $\eta$
instead of $s$, so that every grid is aligned with the data. When a growth
constant is discussed, it is a
set-growth constant $g_S$ as in Definition~\ref{def:growth}, for the
Euclidean distance on $\R^{n_c+n_d}$:
\begin{equation}\label{eq:tu-setgrowth}
F(v)-\OPT\ge g_S\dist(v,\mathcal S)^2\qquad(v\in X),\qquad
\kappa_c=\kappa(\bar L,g_S)=\max\{1,\bar L/g_S\}.
\end{equation}
Point growth at a unique minimizer is the case $\mathcal S=\{v^*\}$. Neither
$g_S$ nor $\mathcal S$ is used by any algorithm below.

\begin{remark}[Directional versus coordinate curvature]\label{rem:tu-curv}
With $\mathcal W=\R^{n_c}$, \eqref{eq:tu-curv} implies, for each
$z\in\mathcal Z$, the coordinate curvature condition of
Definition~\ref{def:curvature} with $L_i=\bar L$ for the continuous
coordinates, but not conversely. For a quadratic, every positive rational
that is at least the largest absolute row sum of $H_{xx}$ is a valid
$\bar L$. For $\mathcal W=\ker C$ and a quadratic, $\bar L$ is valid if and
only if $\Xi_C^{\T}(\bar LI-H_{xx})\Xi_C\succeq0$ for a basis matrix $\Xi_C$
of $\ker C$; no orthonormality of $\Xi_C$ is needed, since $w=\Xi_Cc$ gives
$w^{\T}(\bar LI-H_{xx})w=c^{\T}\Xi_C^{\T}(\bar LI-H_{xx})\Xi_Cc$.
Alternatively, if $C$ has full row rank and
$\Pi_C=I-C^{\T}(CC^{\T})^{-1}C$, every positive rational that is at least
the largest absolute row sum of $\Pi_CH_{xx}\Pi_C$ is valid, because
$\Pi_Cw=w$ for $w\in\ker C$ and the spectral radius of a symmetric matrix is
at most its largest absolute row sum. Both checks are rational and change no
factor scope. Example~\ref{ex:tu-fullcurv} below shows that coordinate
curvature alone is not enough.
\end{remark}

\begin{proposition}[Checking alignment without enumerating $\mathcal Z$]\label{prop:tu-align}
Fix $z^0\in\mathcal Z$ and let $B_k$ be the $k$th column of $B$. Condition
\eqref{eq:tu-align} holds if and only if $\ell/\eta$ and $u/\eta$ are integral,
$(b-Bz^0)/\eta\in\Z^m$, and $B_k(t-z^0_k)/\eta\in\Z^m$ for every $k$ and
every $t\in Z_k$. This takes $1+\sum_k\lvert Z_k\rvert$ vector tests.
\end{proposition}

\begin{proof}
For every $z\in\mathcal Z$,
$(b-Bz)/\eta=(b-Bz^0)/\eta-\sum_kB_k(z_k-z^0_k)/\eta$, a sum of integral
vectors under the stated tests. Conversely, if \eqref{eq:tu-align} holds, the
tests are differences of two vectors $(b-Bz)/\eta$ for points of
$\mathcal Z$ that differ in at most one coordinate.
\end{proof}

\subsection{Feasible rounding on aligned cells}\label{sec:tu-rounding}

\begin{definition}[Level-$j$ domains]\label{def:tu-domain}
A \emph{level-$j$ domain} is a product
$\mathcal D=\mathcal D_1\times\dots\times\mathcal D_{n_c}\times\Lambda_1\times\dots\times\Lambda_{n_d}$,
where each $\mathcal D_i\subseteq[\ell_i,u_i]$ is a nonempty union of finitely
many closed intervals $[a,a']$, $a\le a'$, with $a,a'\in h_j\Z$, and each
$\Lambda_k\subseteq Z_k$ is nonempty. Its grid in coordinate $i$ is
$G_j(\mathcal D_i)=\mathcal D_i\cap h_j\Z$, its \emph{cells} are the intervals
$[t,t+h_j]\subseteq\mathcal D_i$ with $t\in h_j\Z$, and its grid is
$G_j(\mathcal D)=\prod_iG_j(\mathcal D_i)\times\prod_k\Lambda_k$.
\end{definition}

Since $\ell,u\in\eta\Z^{n_c}\subseteq h_j\Z^{n_c}$, the box
$[\ell,u]\times\mathcal Z$ is a level-$0$ domain, and a level-$j$ domain is
also a level-$(j+1)$ domain. Components of $\mathcal D_i$ may be single points.

\begin{lemma}[Feasible corner rounding]\label{lem:tu-round}
Let $\mathcal D$ be a level-$j$ domain and $(x,z)\in X\cap\mathcal D$. There is
a random vector $Y\in\R^{n_c}$ with finitely many values such that, almost
surely, $(Y,z)\in X\cap G_j(\mathcal D)$ and $Y-x\in\mathcal W$, and
\begin{enumerate}[label=(\alph*)]
\item $\E Y=x$;
\item if $x_i\in h_j\Z$ then $Y_i=x_i$; otherwise, with
$t_i=h_j\lfloor x_i/h_j\rfloor$, the cell $[t_i,t_i+h_j]$ lies in $\mathcal D_i$
and $Y_i\in\{t_i,t_i+h_j\}$;
\item $\E\norm{Y-x}^2\le n_ch_j^2/4$.
\end{enumerate}
\end{lemma}

\begin{proof}
Write $h=h_j$ and $t_i=h\lfloor x_i/h\rfloor$, so $x=t+h\vartheta$ with
$\vartheta\in[0,1)^{n_c}$. Let
$\mathcal I_0=\{i:\vartheta_i=0\}=\{i:x_i\in h\Z\}$. If $i\notin\mathcal I_0$,
then $x_i$ lies in a component $[a,a']$ of $\mathcal D_i$ with $a,a'\in h\Z$
and $a<x_i<a'$, so $a\le t_i$ and $t_i+h\le a'$; the cell $[t_i,t_i+h]$ lies
in $\mathcal D_i$.

Consider the polytope
\[
\mathcal Q=\Bigl\{\vartheta'\in\R^{n_c}:\ A\vartheta'\le\frac{b-Bz-At}{h},\
0\le\vartheta'\le\mathbf 1,\ \vartheta'_i\le0\ (i\in\mathcal I_0)\Bigr\}.
\]
The right-hand side $(b-Bz-At)/h$ is integral:
$(b-Bz)/h=2^j(b-Bz)/\eta\in\Z^m$ by \eqref{eq:tu-align}, and $At/h\in\Z^m$
because $A$ is integral and $t/h\in\Z^{n_c}$. Appending unit rows to a TU
matrix gives a TU matrix \cite[Section~19.1]{Schrijver1986}, so by the
Hoffman--Kruskal theorem \cite{HoffmanKruskal1956} every vertex of $\mathcal Q$
is integral, hence lies in $\{0,1\}^{n_c}$ and vanishes on $\mathcal I_0$.
Since $(x,z)\in X$, $\vartheta\in\mathcal Q$, and the polytope $\mathcal Q$ is
the convex hull of its vertices, so $\vartheta=\sum_k\lambda_k\vartheta^k$ with
vertices $\vartheta^k$ and convex weights $\lambda_k$. Let
$Y=t+h\vartheta^k$ with probability $\lambda_k$.

Then $\E Y=t+h\vartheta=x$. For each outcome,
$AY+Bz=At+hA\vartheta^k+Bz\le b$. For $i\in\mathcal I_0$, $Y_i=t_i=x_i$; for
$i\notin\mathcal I_0$, $Y_i\in\{t_i,t_i+h\}$, which lies in the cell
$[t_i,t_i+h]\subseteq\mathcal D_i\subseteq[\ell_i,u_i]$. Hence
$(Y,z)\in X\cap G_j(\mathcal D)$ and (b) holds. If $\mathcal W=\ker C$, each row
$A_r$ of $C$ belongs to an equality pair, so $A_rY+B_rz=b_r=A_rx+B_rz$ and
$A_r(Y-x)=0$. Finally, for $i\notin\mathcal I_0$ the variable $Y_i$ takes two
values with mean $x_i$, so $\E(Y_i-x_i)^2=(x_i-t_i)(t_i+h-x_i)\le h^2/4$, and
coordinates in $\mathcal I_0$ contribute zero. This proves (c).
\end{proof}

The rounded coordinates are dependent: they follow a vertex decomposition of
the polytope $\mathcal Q$, not a product law. This is why the next lemma needs
\eqref{eq:tu-curv} rather than coordinate curvature.

\begin{lemma}[Rounding allowance]\label{lem:tu-allow}
In the situation of Lemma~\ref{lem:tu-round},
\[
\E F(Y,z)\le F(x,z)+\frac{\bar L}2\E\norm{Y-x}^2\le F(x,z)+E_j .
\]
\end{lemma}

\begin{proof}
Fix an outcome and put $w=Y-x\in\mathcal W$. The segment from $x$ to $Y$ lies
in the box $[\ell,u]$, and $\phi(t)=F(x+tw,z)$ satisfies
$\phi''(t)=w^{\T}\nabla^2_{xx}F(x+tw,z)w\le\bar L\norm w^2$ by
\eqref{eq:tu-curv}. Taylor's formula with integral remainder gives
$\phi(1)\le\phi(0)+\phi'(0)+\bar L\norm w^2/2$, that is,
$F(Y,z)\le F(x,z)+\nabla_xF(x,z)^{\T}(Y-x)+\frac{\bar L}2\norm{Y-x}^2$.
Take expectations and use $\E(Y-x)=0$ and Lemma~\ref{lem:tu-round}(c).
\end{proof}

\begin{example}[Coordinate curvature is not enough]\label{ex:tu-fullcurv}
Let $0<h\le1/2$, $X=\{x\in[0,1]^2:x_1=x_2\}$ and
$F(x)=2x_1x_2-h(x_1+x_2)$. Both coordinate curvatures are zero, so the
corrections of Proposition~\ref{prop:cellwise} vanish. On $X$,
$F(t,t)=2t^2-2ht$ has minimum $\OPT=-h^2/2$ at $t=h/2$, whereas every feasible
point of the mesh-$h$ grid has value $F(kh,kh)=2h^2k(k-1)\ge0$. Corrections
computed from coordinate curvature would therefore certify the false bound
$0$. The Hessian has eigenvalues $\pm2$, so $\bar L=2$ is valid with
$\mathcal W=\R^2$, and the allowance $n_c\bar Lh^2/8=h^2/2$ gives the correct
bound $0-h^2/2=\OPT$; the allowance is attained.
\end{example}

\subsection{Filtering and certificates}\label{sec:tu-filter}

The tree computation of Lemma~\ref{lem:dp} applies to $F$ on a grid
$G_j(\mathcal D)$ after each row $r$ is assigned to a bag containing its
support, as the table $\iota_r(y)=0$ if $A_ry_x+B_ry_z\le b_r$ and
$\iota_r(y)=+\infty$ otherwise. The proof of Lemma~\ref{lem:dp} only adds and
minimizes tables, so it holds for values in $\R\cup\{+\infty\}$; when outgoing
messages are formed by subtracting one incoming term from a sum, one stores
the sum of the finite terms and the number of infinite ones. For a level-$j$
domain $\mathcal D$ it returns
\[
\mu=\min\{F(y):y\in X\cap G_j(\mathcal D)\},\qquad
m_i(t)=\min\{F(y):y\in X\cap G_j(\mathcal D),\ y_i=t\}-E_j
\]
for every coordinate $i$ and every node $t$ of that coordinate (for the $k$th
discrete coordinate we write $m_k$), and a minimizing grid point, with $\min\emptyset=+\infty$, using
$O\bigl(p(\lvert\mathcal A\rvert+n+N+m)K^p\bigr)$ table operations, where $K$
bounds the number of nodes per coordinate. The correction is the constant
$E_j$, so $m_i$ is the min-marginal of $Q=F-E_j$ over the feasible grid
points, in the notation of Definition~\ref{def:corr}, and $\mu-E_j$ is the
minimum of $Q$.

\begin{algorithm}[TU-GRID: grids under TU coupling]\label{alg:tugrid}
Input: the data of Definition~\ref{def:tu-model} and an accuracy
$\varepsilon>0$. Set $\mathcal D^{(0)}=[\ell,u]\times\mathcal Z$. For
$j=0,1,2,\dots$:
\begin{enumerate}[label=(\arabic*)]
\item Compute $\mu_j$, a minimizer $y^{(j)}\in X\cap G_j(\mathcal D^{(j)})$, and
the min-marginals $m_i$ on $G_j(\mathcal D^{(j)})$. If $\mu_0=+\infty$, stop and
report $X=\emptyset$.
\item Put $U_j=\mu_j$ and $\beta_j=\mu_j-E_j$. If $E_j\le\varepsilon$, return
$y^{(j)}$ and $[\beta_j,U_j]$.
\item (Union filter.) For each continuous $i$, let $\mathcal D^{(j+1)}_i$ be
the union of all cells $[t,t+h_j]$ of $\mathcal D^{(j)}_i$ with
$\min\{m_i(t),m_i(t+h_j)\}\le U_j$ and of all single-point components
$\{t\}$ of $\mathcal D^{(j)}_i$ with $m_i(t)\le U_j$. For each $k$, let
$\Lambda^{(j+1)}_k=\{t\in\Lambda^{(j)}_k:m_{k}(t)\le U_j\}$.
\end{enumerate}
The \emph{hull variant} replaces each $\mathcal D^{(j+1)}_i$ by its convex hull.
\end{algorithm}

\begin{proposition}[Soundness]\label{prop:tu-sound}
For TU-GRID and for its hull variant, for every level $j$ that is reached:
\begin{enumerate}[label=(\alph*)]
\item if $X\ne\emptyset$ then $\mu_0<\infty$;
\item $\beta_j\le F(v)$ for every $v\in X\cap\mathcal D^{(j)}$;
\item every $v\in(X\cap\mathcal D^{(j)})\setminus\mathcal D^{(j+1)}$ has
$F(v)>U_j$;
\item $\mathcal D^{(j+1)}$ is a level-$(j+1)$ domain that contains $\mathcal S$
and $y^{(j)}$; hence $y^{(j)}\in X\cap G_{j+1}(\mathcal D^{(j+1)})$ and
$U_{j+1}\le U_j$;
\item $\beta_j\le\OPT\le U_j\le\OPT+E_j$ and $U_j-\beta_j=E_j$.
\end{enumerate}
Moreover, a checker that recomputes every $\mu_j$ and $m_i$ from the stored
domains and verifies every filtering decision obtains $\beta_J\le\OPT$ for
the last level $J$ from (b) and (c) alone, without using any minimizer,
uniqueness or growth.
\end{proposition}

\begin{proof}
(a) Apply Lemma~\ref{lem:tu-round} with $j=0$ to any $(x,z)\in X$; its
outcomes are feasible points of $G_0(\mathcal D^{(0)})$.

(b) Let $v=(x,z)\in X\cap\mathcal D^{(j)}$ and let $Y$ be as in
Lemma~\ref{lem:tu-round}. Every outcome $(Y,z)$ is a feasible grid point, so
$\mu_j\le\E F(Y,z)\le F(x,z)+E_j$ by Lemma~\ref{lem:tu-allow}.

(c) Let $v=(x,z)\in(X\cap\mathcal D^{(j)})\setminus\mathcal D^{(j+1)}$ and
let $Y$ be as in Lemma~\ref{lem:tu-round}. If some value
$z_k\notin\Lambda^{(j+1)}_k$, every outcome has the same value $z_k$, so
$m_k(z_k)\le F(Y,z)-E_j$ almost surely and $m_k(z_k)\le\E F(Y,z)-E_j\le F(v)$;
removal means $m_k(z_k)>U_j$. Otherwise some
$x_i\notin\mathcal D^{(j+1)}_i$. If $x_i\in h_j\Z$, then $Y_i=x_i$ almost
surely and as before $m_i(x_i)\le F(v)$. Since $x_i\in\mathcal D^{(j)}_i$,
either $\{x_i\}$ is a single-point component, which was removed, so
$m_i(x_i)>U_j$; or $x_i$ is an endpoint of at least one cell, all such cells
were removed, and each removed cell has both endpoint values above $U_j$, so
again $m_i(x_i)>U_j$. If $x_i\notin h_j\Z$, then $Y_i\in\{t_i,t_i+h_j\}$ for
the cell $[t_i,t_i+h_j]\subseteq\mathcal D^{(j)}_i$ containing $x_i$, which was
removed, and $F(v)\ge\E F(Y,z)-E_j\ge\min\{m_i(t_i),m_i(t_i+h_j)\}>U_j$. In
the hull variant the retained set is larger, so (c) holds a fortiori.

(d) The retained sets are finite unions of closed intervals with endpoints in
$h_j\Z\subseteq h_{j+1}\Z$, as are their hulls. $X$ is compact and $F$ is
continuous, so $\mathcal S\ne\emptyset$ when $X\ne\emptyset$. By induction
$\mathcal S\subseteq\mathcal D^{(j)}$; minimizers have value
$\OPT\le F(y^{(j)})=U_j$, so by (c) they are not removed. The point $y^{(j)}$
has value $U_j$, so it is not removed either; its continuous coordinates lie
in $h_j\Z\subseteq h_{j+1}\Z$, hence $y^{(j)}\in X\cap G_{j+1}(\mathcal D^{(j+1)})$
and $\mu_{j+1}\le\mu_j$. In particular every $\mathcal D^{(j+1)}_i$ and
$\Lambda^{(j+1)}_k$ is nonempty.

(e) $\OPT\le U_j$ because $y^{(j)}\in X$. By (d) a minimizer lies in
$\mathcal D^{(j)}$, so (b) gives $\beta_j\le\OPT$; rounding that minimizer as
in (b) gives $\mu_j\le\OPT+E_j$.

For the last claim, let $v\in X$. If $v\in\mathcal D^{(J)}$, then
$F(v)\ge\beta_J$ by (b). Otherwise $v$ is removed at a first level $j<J$, and
$F(v)>U_j\ge U_J\ge\beta_J$ by (c); the inequality $U_j\ge U_J$ is checked
directly from the recomputed values. Hence $\OPT\ge\beta_J$.
\end{proof}

The certificate record consists of the data, a TU certificate for $A$
(a structural certificate such as a network or interval representation, or a
run of a polynomial-time TU recognition algorithm \cite[Chapter~20]{Schrijver1986}),
the curvature certificate for \eqref{eq:tu-curv}, $\eta$, the domains, the
tables, and the points $y^{(j)}$. As in Theorem~\ref{thm:certificate}, the
whole history is needed: the last grid bounds $F$ only on
$X\cap\mathcal D^{(J)}$.

\subsection{Accuracy-independent node counts}\label{sec:tu-states}

\begin{theorem}[Node counts under set growth]\label{thm:tu-states}
Assume \eqref{eq:tu-setgrowth} and put $a_j=\sqrt{n_c\bar L/(4g_S)}\,h_j$. For
every level $j$ reached by TU-GRID:
\begin{enumerate}[label=(\alph*)]
\item every point of $\mathcal D^{(j+1)}_i$ lies within $a_j+h_j$ of
$\mathcal S_i$, and every value in $\Lambda^{(j+1)}_k$ lies within $a_j$ of
$\mathcal S_k$;
\item if a continuous projection $\mathcal S_i$ has at most $r_i$ points, then
\[
\bigl\lvert G_{j+1}(\mathcal D^{(j+1)}_i)\bigr\rvert\le
r_i\Bigl(5+\bigl\lfloor2\sqrt{n_c\bar L/g_S}\bigr\rfloor\Bigr)
\le r_i\bigl(5+\bigl\lfloor2\sqrt{n_c\kappa_c}\bigr\rfloor\bigr);
\]
\item in the hull variant the statement on values in (a) holds; if moreover
$\mathcal S=\{v^*\}$, every point of $\mathcal D^{(j+1)}_i$ lies within
$a_j+h_j$ of $v^*_i$, and the bound in (b) holds with $r_i=1$.
\end{enumerate}
\end{theorem}

\begin{proof}
(a) A retained cell has an endpoint $t$ with $m_i(t)\le U_j$, a retained
single point $\{t\}$ has $m_i(t)\le U_j$, and a retained value $t$ has
$m_k(t)\le U_j$. In each case $m_i(t)<\infty$, so there is a feasible grid
point $w$ with $w_i=t$ and, by Proposition~\ref{prop:tu-sound}(e),
\[
F(w)=m_i(t)+E_j\le U_j+E_j\le\OPT+2E_j .
\]
By \eqref{eq:tu-setgrowth}, $\dist(w,\mathcal S)^2\le2E_j/g_S=a_j^2$, so some
$v\in\mathcal S$ has $\lvert t-v_i\rvert\le\norm{w-v}\le a_j$. Points of a cell
are within $h_j$ of each of its endpoints.

(b) By (a), $\mathcal D^{(j+1)}_i$ lies in the union over $\nu\in\mathcal S_i$
of the intervals $[\nu-a_j-h_j,\nu+a_j+h_j]$. Each has length
$2a_j+2h_j=h_{j+1}(4a_j/h_j+4)$, so it contains at most
$\lfloor4a_j/h_j\rfloor+5$ points of $h_{j+1}\Z$, and
$4a_j/h_j=2\sqrt{n_c\bar L/g_S}$. Finally $\bar L/g_S\le\kappa_c$.

(c) In the hull variant the argument of (a) applies to the retained cells,
single points and values before the hull is taken. If $\mathcal S=\{v^*\}$,
the retained cells and points lie in the interval
$[v^*_i-a_j-h_j,v^*_i+a_j+h_j]$, and hence so does their hull.
\end{proof}

The bound is per optimal coordinate value: different coordinates may be
near different minimizers, and no optimal combination is enumerated;
$\mathcal S$ can have up to $\prod_ir_i\prod_k\lvert Z_k\rvert$ points. The
finite-projection hypothesis cannot be dropped for TU-GRID: if $\mathcal S_i$
contains an interval $[a,a']$ with $a<a'$, then
$[a,a']\subseteq\mathcal D^{(j)}_i$ for every $j$ by
Proposition~\ref{prop:tu-sound}(d), so the level-$j$ grid has at least
$(a'-a)/h_j-1$ nodes in coordinate $i$, which grows with the accuracy.

\begin{example}[Two minimizers per block: union versus hull]\label{ex:tu-union}
For one block take continuous $x_1,x_2,x_3\in[0,1]$ with the TU equality
$x_1+x_2=x_3$ and
\[
F_b(x)=(2x_1-x_3)^2+\tfrac14x_3(1-x_3).
\]
Both terms are nonnegative on the feasible set $X_b$ of the block, so
$\mathcal S_b=\{(0,0,0),(\tfrac12,\tfrac12,1)\}$ and $\OPT=0$. On $X_b$ put
$w=x_1-x_3/2$, so $x_2=x_3/2-w$ and $F_b=4w^2+x_3(1-x_3)/4$. If $x_3\le1/2$,
the squared distance to $(0,0,0)$ is $2w^2+\frac32x_3^2$ and
$x_3(1-x_3)/4\ge x_3/8\ge x_3^2/4$, so
$F_b\ge4w^2+x_3^2/4\ge\frac16(2w^2+\frac32x_3^2)$. If $x_3\ge1/2$, the same
computation with $1-x_3$ in place of $x_3$ and the minimizer
$(\frac12,\frac12,1)$ gives the same inequality. Hence $g_S=1/6$. The Hessian
has rows $(8,0,-4)$, $(0,0,0)$, $(-4,0,3/2)$, so $\bar L=12$ is valid. For $n_b$
independent blocks, $n_c=3n_b$, and the bags are the blocks ($p=3$). The
optimal set is the product of the block sets, so squared distances to it add
over blocks and $g_S=1/6$ and $\bar L=12$ are unchanged; every projection has
$r_i=2$ points, and there are $2^{n_b}$ minimizers. With $\eta=1$, the hull
variant keeps $[0,1]$ in the coordinate $x_3$ of every block at every level,
because both $0$ and $1$ are optimal values, so its level-$j$ grid has
$2^j+1$ nodes there.
The union variant has at most $2(5+\lfloor2\sqrt{216n_b}\rfloor)$ nodes per
coordinate at every level $j\ge1$ by Theorem~\ref{thm:tu-states}.
\end{example}

\subsection{Complexity of certified approximation}\label{sec:tu-complexity}

\begin{theorem}[Certified approximation with TU coupling]\label{thm:tu-approx}
Assume Definition~\ref{def:tu-model} with factors that are explicit rational
quadratics or explicitly encoded rational polynomials of fixed degree, and
let $I$ be the binary input length, including $\eta$, $\bar L$ and the
decomposition. For $\varepsilon>0$, TU-GRID$(\varepsilon)$ either certifies
$X=\emptyset$ at level $0$, or stops at
\[
J=\min\{j\ge0:E_j\le\varepsilon\}\le
\max\Bigl\{0,\Bigl\lceil\tfrac12\log_2\frac{n_c\bar L\eta^2}{8\varepsilon}\Bigr\rceil\Bigr\}
\]
with a feasible rational point and an interval $[\beta_J,U_J]\ni\OPT$ of width
$E_J\le\varepsilon$, certified as in Proposition~\ref{prop:tu-sound}. All
numbers have bit length polynomial in $I+J$, and the number of table
operations is
\[
O\Bigl(p\,(\lvert\mathcal A\rvert+n+N+m)\sum_{j=0}^JK_j^p\Bigr),\qquad
K_0\le\max\{K_Z,s/\eta+1\},
\]
where $K_j$ is the largest number of nodes per coordinate at level $j$.
Under \eqref{eq:tu-setgrowth} with $\lvert\mathcal S_i\rvert\le r$ for every
continuous $i$, $K_j\le K=\max\{K_Z,r(5+\lfloor2\sqrt{n_c\kappa_c}\rfloor)\}$
for every $j\ge1$. The algorithm uses neither $g_S$ nor $r$.
\end{theorem}

\begin{proof}
Correctness is Proposition~\ref{prop:tu-sound}, and the width is
$U_J-\beta_J=E_J$. The bound on $J$ follows from $E_j=n_c\bar L\eta^24^{-j}/8$.
The table count is the tree computation above, summed over levels; the
level-$0$ grid has $(u_i-\ell_i)/\eta+1\le s/\eta+1$ nodes in a continuous
coordinate and at most $K_Z$ values in a discrete one, and values are never
added. The bound $K$ is Theorem~\ref{thm:tu-states}(b) together with
$\lvert\Lambda^{(j)}_k\rvert\le K_Z$. Grid points at level $j$ lie in
$h_j\Z\cap[\ell,u]$, so their coordinates are rationals with denominators
dividing $\Delta_\eta2^j$ and numerators bounded by
$\Delta_\eta2^j\max_i\max\{\lvert\ell_i\rvert,\lvert u_i\rvert\}$. Explicit
fixed-degree factors evaluated at such points have polynomial bit length,
and every table entry is a sum of at most $\lvert\mathcal A\rvert$ factor
values or a minimum of such sums.
\end{proof}

Without growth the same certificate is produced, but the node counts depend
on the accuracy: a continuous coordinate has at most $s2^j/\eta+1$ nodes at
level $j$, and if $J\ge1$, minimality of $J$ gives
$2^J<\eta\sqrt{n_c\bar L/(2\varepsilon)}$. A bag of $p$ continuous
coordinates then has fewer than $(s\sqrt{n_c\bar L/(2\varepsilon)}+1)^p$
table entries at the last level, which is polynomial in $1/\varepsilon$, not
in $\log(1/\varepsilon)$.

Under set growth with finitely many optimal values per coordinate, TU-GRID
runs in time polynomial in $I+\log(1/\varepsilon)$ for fixed $p$ when the
condition number $\kappa_c$, the box widths measured in mesh units
($s/\eta$), and $r$ are polynomially bounded; recall that $K_Z\le I$. It is
not fixed-parameter tractable: its running time is not bounded by
$f(p,\kappa_c)\poly(I+\log(1/\varepsilon))$, because $K^p$ contains the
factor $n_c^{p/2}$ and Section~\ref{sec:tu-limits} shows that this factor is
incurred. The level-$0$ grid has up to $s/\eta+1$ nodes per coordinate, which
is pseudopolynomial in the data. A larger common unit $\eta$, checked by
Proposition~\ref{prop:tu-align}, makes this smaller, but unless
$\mathrm P=\mathrm{NP}$ the dependence on $s/\eta$ cannot be removed.
Indeed, Proposition~\ref{lim:prop:constraints}
(Section~\ref{sec:limits-constraints}) gives instances of
Definition~\ref{def:tu-model} with $p=3$, a unique minimizer and $K_Z=2$ on
which approximating $\OPT$ within $\frac12$ is NP-hard. Their objective is
affine, so every $\bar L>0$ is valid, and $\bar L\le g_S$ gives
$\kappa_c=1$; of the quantities above, only $s/\eta$ is not polynomially
bounded on them.

\subsection{Exact output for quadratic objectives}\label{sec:tu-exact}

In this subsection $F(v)=\frac12v^{\T}Hv+\nabla F(0)^{\T}v+F(0)$ is a
quadratic with rational data on the model of Definition~\ref{def:tu-model}.
Let $\Delta$ be a positive common denominator of its monomial coefficients
$H_{ii}/2$, $H_{ik}$ $(i<k)$, the entries of $\nabla F(0)$, and $F(0)$, and
put $\hat H=\Delta H$, which is
integral and symmetric, with continuous block $\hat H_{xx}$ and mixed block
$\hat H_{xz}$. For $z\in\mathcal Z$ the continuous part of $X$ is the polytope
\[
\mathcal P_z=\{x\in\R^{n_c}:\hat Ax\le\hat b_z\},\qquad
\hat A=\begin{pmatrix}A\\ I\\ -I\end{pmatrix}\in\{0,\pm1\}^{m'\times n_c},\qquad
\hat b_z=\begin{pmatrix}b-Bz\\ u\\ -\ell\end{pmatrix}\in\Delta_\eta^{-1}\Z^{m'},
\]
with $m'=m+2n_c$, so that $\{x:(x,z)\in X\}=\mathcal P_z$. Put
$c_z=\hat H_{xz}z+\Delta\nabla_xF(0)\in\Z^{n_c}$, so that
$\Delta\nabla_xF(x,z)=\hat H_{xx}x+c_z$. The \emph{TU height constants} are
\begin{equation}\label{eq:tu-constants}
C_{\hat H}=\max\Bigl\{1,\max_{i,k\le n_c}\lvert(\hat H_{xx})_{ik}\rvert\Bigr\},\quad
R_{\mathrm{TU}}=(2n_cC_{\hat H})^{n_c},\quad
\Omega_{\mathrm{TU}}=\Delta(\Delta_\eta R_{\mathrm{TU}})^2,\quad
\tau_{\mathrm{TU}}=\frac1{4m'\Delta_\eta R_{\mathrm{TU}}} .
\end{equation}
They are computed from the original data and have bit length polynomial in
$I$. They take the place of the constants $R$, $\Omega$ and $\tau$ of
Section~\ref{sec:exact}, where all rows are bound rows and the stationarity
systems are principal submatrices of the Hessian; with TU rows a Hadamard
bound on the rows of a saddle-point matrix replaces the diagonal bound.
Appendix~\ref{app:tu} proves the analogues of Lemma~\ref{lem:statpoly} and
Corollary~\ref{cor:height}: some global minimizer has continuous coordinates
in $\rho^{-1}\Z$ for an integer $1\le\rho\le\Delta_\eta R_{\mathrm{TU}}$, every
isolated global minimizer has this property, and $\OPT=a/W$ in lowest terms
with $1\le W\le\Omega_{\mathrm{TU}}$ (Corollary~\ref{cor:tu-height}).
Proposition~\ref{prop:accept} therefore applies with
$\Omega'=\Omega_{\mathrm{TU}}$: if $\beta\le\OPT$, $v\in X$ and $F(v)=a/W$ in
lowest terms with $F(v)-\beta<1/(\Omega_{\mathrm{TU}}W)$, then $v$ is a global
minimizer.

Recovery snaps rows instead of coordinates: the rows of $\hat A$ that are
nearly active at the grid minimizer are made active, and stationarity along
the remaining face is imposed by linear equations.

\begin{algorithm}[TU-EXACT: exact output with TU coupling]\label{alg:tuexact}
Run TU-GRID (Algorithm~\ref{alg:tugrid}) without the stopping test of
step~(2). Between steps~(2) and~(3) of every level $j$, let $(y,z)=y^{(j)}$
and:
\begin{enumerate}[label=(\roman*)]
\item let $\mathcal J=\{r:(\hat b_z)_r-\hat A_ry\le\tau_{\mathrm{TU}}\}$, and
compute a basis matrix $\Xi_{\mathcal J}$ of $\ker\hat A_{\mathcal J}$ by
Gaussian elimination, where $\hat A_{\mathcal J}$ consists of the rows of
$\hat A$ in $\mathcal J$;
\item find by exact linear programming a point $\hat x$ of
\[
\mathcal P_z(\mathcal J)=\bigl\{x\in\mathcal P_z:\
\hat A_{\mathcal J}x=(\hat b_z)_{\mathcal J},\
\Xi_{\mathcal J}^{\T}(\hat H_{xx}x+c_z)=0\bigr\},
\]
or continue with step~(3) if $\mathcal P_z(\mathcal J)$ is empty;
\item if $F(\hat x,z)-\beta_j<1/(\Omega_{\mathrm{TU}}W)$, where $W$ is the
reduced denominator of $F(\hat x,z)$, return $(\hat x,z)$, $F(\hat x,z)$ and
the certificate of levels $0,\dots,j$; otherwise continue with step~(3).
\end{enumerate}
\end{algorithm}

The large denominators of $y$ enter only the comparisons that define
$\mathcal J$. The rest of steps~(i) and~(ii) uses only the original data,
the index set $\mathcal J$ and the integer vector $z$, so it takes time
polynomial in $I$ \cite[Theorem~6.4.12]{GrotschelLovaszSchrijver1988}.

\begin{theorem}[Exact output with TU coupling]\label{thm:tu-exact}
Let $F$ be a rational quadratic on a nonempty set $X$ as in
Definition~\ref{def:tu-model}.
\begin{enumerate}[label=(\alph*)]
\item Every point returned by TU-EXACT is a global minimizer, and its value
is $\OPT$. This uses neither uniqueness nor growth.
\item TU-EXACT terminates.
\item Under \eqref{eq:tu-setgrowth} it terminates by the first level $j$ with
$E_j\le\min\{g_S\tau_{\mathrm{TU}}^2/(4n_c),\,1/(2\Omega_{\mathrm{TU}}^2)\}$,
hence after at most
\[
J_{\rm ex}=\max\Bigl\{0,\Bigl\lceil\log_2\frac{\eta n_c\sqrt{\kappa_c}}{\tau_{\mathrm{TU}}}\Bigr\rceil,
\Bigl\lceil\log_2\bigl(\eta\Omega_{\mathrm{TU}}\sqrt{n_c\bar L}\bigr)\Bigr\rceil\Bigr\}
=\poly(I)+O(\log\kappa_c)
\]
levels. If moreover every continuous projection $\mathcal S_i$ has at most
$r$ points, the table operations are as in Theorem~\ref{thm:tu-approx} with
$J$ replaced by $J_{\rm ex}$, and steps~(i)--(iii) add a cost polynomial in
$I+j$ at level $j$.
\end{enumerate}
\end{theorem}

The proof is in Appendix~\ref{app:tu}. It follows the proofs of
Theorems~\ref{thm:transfer} and~\ref{thm:exact}, with a snapping lemma for
rows of $\hat A$ (Lemma~\ref{lem:tu-snap}) in place of
Lemma~\ref{lem:snap}.

\subsection{The uniform mesh: cost and non-uniform alternatives}\label{sec:tu-limits}

The running time of CT (Algorithm~\ref{alg:ct}) is bounded by
$f(p,\bar\kappa)\poly(I+q)$ because its graded grids cover the localization
radius with $O(\sqrt{\bar\kappa}\log(n+2))$ nodes per coordinate
(Section~\ref{sec:growth}). With TU coupling the radius is
$a_j+h_j$, of order $\sqrt{n_c\kappa_c}\,h_j$ (Theorem~\ref{thm:tu-states}),
but the mesh is uniform. The factor $n_c^{p/2}$ in $K^p$ is incurred:
Corollary~\ref{cor:uniformgrid} shows that filtered uniform grids keep at
least $\sqrt{(n-2)\kappa}+1$ nodes per coordinate on the box instance of
Proposition~\ref{prop:sharp}, and its proof applies to TU-GRID, whose
allowance $E_j$ is at least the sum of the coordinate corrections of a
uniform grid. With a redundant TU row that forces a bag of size $p$,
TU-GRID$(\varepsilon)$ then builds a table with at least $n_c^{p/2}$ entries
for every $\varepsilon\le\bar L/64$, while $\kappa_c=2$
(Lemma~\ref{lem:tu-uniform} in Appendix~\ref{app:tu}).

Natural non-uniform alternatives fail, as follows.

\begin{proposition}[Misaligned grids give invalid bounds]\label{prop:tu-misaligned}
Let $G_1,G_2\subseteq[0,1]$ be finite sets containing $0$ and $1$, let
$d_i:G_i\to[0,\infty)$ be arbitrary, let $L>0$, and let $m\in(0,1)$ satisfy
\begin{equation}\label{eq:tu-gap}
L(\nu-m)^2>d_1(\nu)+d_2(\nu)\qquad\text{for every }\nu\in G_1\cap G_2 .
\end{equation}
For $\lambda\ge0$ let $X=\{x\in[0,1]^2:x_1-x_2\le0\}$ and
$F_\lambda(x)=\frac L2\bigl((x_1-m)^2+(x_2-m)^2\bigr)+\lambda(x_2-x_1)$. Then
$\nabla^2F_\lambda=LI$, $\OPT=0$ is attained only at $(m,m)$, and
$F_\lambda(x)\ge\frac L2\norm{x-(m,m)}^2$ on $X$, for every $\lambda$. Let
$\delta=\min\{y_2-y_1:y\in G_1\times G_2,\ y_2>y_1\}$. If
$\lambda>(\max d_1+\max d_2)/\delta$, then
\[
\min\bigl\{F_\lambda(y)-d_1(y_1)-d_2(y_2):\ y\in(G_1\times G_2)\cap X\bigr\}>\OPT .
\]
\end{proposition}

\begin{proof}
On $X$, $\lambda(x_2-x_1)\ge0$, which gives the growth inequality and
$\OPT=0$ at $(m,m)$ only. A feasible grid point either has $y_1=y_2=\nu\in
G_1\cap G_2$, with corrected value $L(\nu-m)^2-d_1(\nu)-d_2(\nu)>0$ by
\eqref{eq:tu-gap}, or has $y_2-y_1\ge\delta$, with corrected value at least
$\lambda\delta-\max d_1-\max d_2>0$.
\end{proof}

For instance, $G_1=\{0,\frac12,1\}$, $G_2=\{0,\frac13,1\}$, $m=\frac25$ and the
corrections $d_i(\nu)=\frac L8w_i(\nu)^2$ of Section~\ref{sec:grids} satisfy
\eqref{eq:tu-gap}, since $G_1\cap G_2=\{0,1\}$,
$\frac4{25}>\frac1{32}+\frac1{72}$ and $\frac9{25}>\frac1{32}+\frac1{18}$; here
$\delta=\frac13$, and every $\lambda>\frac{75}{288}L$ yields a corrected grid
minimum above $\OPT$. Graded grids of Section~\ref{sec:growth} built around
two different feasible centers behave in the same way: for $h=\frac1{16}$,
$\theta=\frac14$ and the centers $\frac3{10}$ and $\frac7{10}$ on $[0,1]$, the
two grids have eleven nodes each and only $0$ and $1$ in common, and
\eqref{eq:tu-gap} holds with $m=\frac12$ (Appendix~\ref{app:tu-limits}). The
multiplier $\lambda$ of the order constraint is invisible to the curvature
$L$ and to the growth constant. A correction that depends only on these and
on the grids, and is small enough at the common nodes to satisfy
\eqref{eq:tu-gap}, therefore cannot absorb it. Aligned uniform grids escape
because then $G_1\cap G_2$ contains a node within $h/2$ of $m$, and
$Lh^2/4\le d_1+d_2$ there.

\begin{example}[A graded grid fails on a sum equality]\label{ex:tu-sum}
Let $X=\{x\in[0,3]^3:x_1+x_2-x_3=0\}$, which is TU with integral data, and
$F(x)=\frac L2\norm{x-(1,1,2)}^2$, so $\OPT=0$ at $(1,1,2)\in X$,
$\nabla^2F=LI$ and growth holds with $g=L/2$. Build the graded grid of
Section~\ref{sec:growth} in each coordinate around the feasible center $0$,
with $h=1$ and $\theta=\frac14$: the steps $h+\theta t$ give the common grid
$G=\{0,1,\frac94,3\}$, clipped at $3$, and the corrections
$d(0)=\frac L8$, $d(1)=d(\frac94)=\frac{25L}{128}$, $d(3)=\frac{9L}{128}$. The
corrected grid minimum over the seven feasible grid points is
$\frac{31}{64}L>\OPT$ (Appendix~\ref{app:tu-limits}): with a sum row, a common
graded pattern around a feasible center already gives an invalid bound,
without any multiplier effect. The reason is visible in the cell
$[0,1]\times[0,1]\times[1,\frac94]$: its intersection with the plane
$x_1+x_2=x_3$ is the triangle with vertices $(1,0,1)$, $(0,1,1)$ and
$(1,1,2)$, and the last vertex is not a grid point, so the minimizer is not a
convex combination of feasible grid points of its cell.
\end{example}

\begin{remark}[Comparison with Bienstock and Mu\~noz]\label{rem:tu-bm}
Bienstock and Mu\~noz \cite[Theorems~4 and~15]{BienstockMunoz2018} construct,
for mixed-binary polynomial problems with variables scaled to $[0,1]$, a linear
objective $c$, constraints of degree at most $\pi$ and coefficient $1$-norm at
most $\mathfrak F$, and a constraint intersection graph of treewidth $\omega$, a
linear program with $O\bigl((2\pi/\epsilon)^{\omega+1}n\log(\pi/\epsilon)\bigr)$
variables and constraints whose solutions violate each constraint by at most
$\mathfrak F\epsilon$ and are $\lVert c\rVert_1\epsilon$-optimal. They also argue that,
unless $\mathrm P=\mathrm{NP}$, neither the pseudopolynomial dependence on
$1/\epsilon$ nor the coefficient scaling of the violation can be removed for that
class. The results here differ in four respects. (i) Feasibility is exact, but
only for TU continuous columns with aligned data; their constraints are
arbitrary polynomials. (ii) A nonlinear objective is handled directly through
$\bar L$, whereas their objective is linear and nonlinear objectives are
encoded by additional constraints, which may increase the width. (iii) Under
growth with finitely many optimal coordinate values, the per-level tables of
TU-GRID do not depend on its accuracy $\varepsilon$, and $\varepsilon$ enters
only through $O(\log(1/\varepsilon))$ levels; this does not conflict with
their lower bounds, which concern their whole class and hence instances
without growth or TU structure, and the factor $(n_c\kappa_c)^{p/2}$ can
itself be exponential in $I$ because $\kappa_c$ can be. (iv) Without growth,
if $J\ge1$, a last-level table of TU-GRID for a bag of $p$ continuous
coordinates has fewer than $(s\sqrt{n_c\bar L/(2\varepsilon)}+1)^p$ entries
(Section~\ref{sec:tu-complexity}), an exponent $p/2$ in $1/\varepsilon$
against their $\omega+1$ in $1/\epsilon$ for comparable decompositions; the
halving comes from second-order rounding error with exact feasibility, and
the comparison is only indicative because the tolerances are normalized
differently. Neither result contains the other.
\end{remark}
